\section{Formal verification}\label{sec:ver}
\begingroup\sloppy

The Lean~4 development \cite{LeanRepo,Lean4} uses the toolchain \texttt{v4.34.1}, Mathlib \cite{Mathlib2020} at
the tag of the same name, and the unmodified Lean library accompanying \cite{OpenAI2026}. That library contributes lemmas and a few generic
differential-geometric operators (exterior derivative, chart differential); the definitions of the Hamiltonian and of
the algorithms belong to the development, and none of the library's packing constructions or existence results is
used. The entry
point is five short files in \texttt{verification/}: \texttt{Problem.lean} defines the problem using Mathlib alone;
\texttt{Algorithm.lean} exposes the top level of the construction; \texttt{Statement.lean} states correctness of the
construction; \texttt{Computability.lean} and \texttt{EpochComputability.lean} state correctness and termination of the
finite-output algorithms. Theorems in \texttt{verification/} have no namespace; other names are relative to
\texttt{SymplecticPacking}.

\begin{table}[ht]
\centering
\footnotesize
\renewcommand{\arraystretch}{1.15}
\begin{tabular}{@{}L{0.2\textwidth}L{0.74\textwidth}@{}}
\toprule
Statement & Lean \\
\midrule
Theorem~\ref{thm:intro-4d} & \texttt{four\_dim\_theorem}, with \texttt{General.epochRadiusQ},
 \texttt{General.epochHamiltonianQ} \\
Proposition~\ref{thm:intro-stability} & \texttt{four\_dim\_stability} \\
Theorem~\ref{thm:main}(i) & \texttt{main\_theorem}, \texttt{main\_time\_support}, \texttt{main\_support\_radius}; for the
 flow, \texttt{flow\_is\_integral\_curve}, \texttt{flow\_symplectic\_diffeomorphism}, \texttt{flow\_images\_disjoint} \\
Theorem~\ref{thm:main}(ii) & \texttt{\#check\_explicit General.packingHamiltonian} (a metaprogram, see below) \\
Theorem~\ref{thm:main}(iii) & \texttt{original\_name\_algorithm}, with \texttt{\#check\_constructive} on the machine;
 finite form: \texttt{General.Compile.main\_finite} \\
Corollaries~\ref{thm:main-sy}--\ref{thm:main-cone4} & \texttt{neighbourhood\_criterion},
 \texttt{flow\_neighbourhood\_embedding}; \texttt{equal\_balls}; \texttt{epoch\_packing}, \texttt{epoch\_theorem\_ten},
 \texttt{epoch\_name\_algorithm}, \texttt{epoch\_name\_algorithm\_packing}; \texttt{unequal\_cone} \\
embeddings, by case & $n=2$: \texttt{Branch2.embeddingsEq\_spec} (linear datum), \texttt{Branch2.embeddingsGen\_spec}
 (scalar comparator); $n\ge3$: \texttt{BranchSmall.embeddings\_spec} (all $c_i<\frac C2$),
 \texttt{BranchLarge3.embeddings\_spec} ($n=3$, one $c_j\ge\frac C2$), \texttt{BranchLargeHigh.embeddings\_spec}
 ($n\ge4$, one $c_j\ge\frac C2$) \\
\bottomrule
\end{tabular}
\caption{Formal counterparts of the main statements. Every case of \Sref{Construction}{cons:ham-branch} is covered
by \texttt{main\_theorem} and by the same algorithm theorems.}
\label{tab:ver-coverage}\label{tab:ver-levels}
\end{table}

\subsection*{The statement}

Phase space is $\texttt{Problem.E}~n=\texttt{EuclideanSpace}\ \R\ (\texttt{Fin}\ (2n))$, coordinates $j<n$ and $n+j$
being $x_{j+1}$ and $y_{j+1}$ (named $q$ and $p$ in the file); \texttt{omega0} and \texttt{hamVF} give $\omega_0=\sum_jdx_j\wedge dy_j$ and
$X_H=(\partial_yH,-\partial_xH)$, so $\iota_{X_H}\omega_0=dH$. \texttt{capBall}~$s$~$c=\{z:\pi\|z-s\|^2\le c\}$ and
\texttt{capBallOpen} is its strict version. \texttt{Admissible}, \texttt{BallsDisjoint} and \texttt{PackingSpec} are
Definitions~\ref{def:intro-admissible} and~\ref{def:intro-spec} and the disjointness hypothesis, field by field;
\texttt{PackingSpec} quantifies over every integral curve on $[0,1]$, so it does not depend on a chosen solution.
\texttt{EpochSpecR} is the problem as posed, with the closed Euclidean unit balls about
\texttt{epochCenter}~$i=(3i,0,0,0)$, $0\le i<k$ (the centre $s_{i+1}$), and the closed target of radius $R$.

Each theorem is a statement about the defined term. \texttt{Algorithm.lean} shows by \texttt{rfl} that
\texttt{General.packingHamiltonian} is \texttt{HamEnd.hamiltonianOfEmbeddings} applied to
\texttt{General.selectEmbeddings}, the case distinction of \Sref{Construction}{cons:ham-branch}, and that the Epoch
Hamiltonian is its instance with $C=\pi R^2$. The Lean statements agree with the paper, with these differences of
form: \texttt{equal\_balls} does not assume $\eps>0$; \texttt{four\_dim\_stability} also records $R<R'$;
\texttt{four\_dim\_theorem} contains both \texttt{EpochSpecR} and the packing specification; the algorithm theorems
assert \texttt{Nat.RecursiveIn} relative to the oracles given by the input names, and that one program serves every
input is certified by \texttt{\#check\_constructive}; \texttt{flow\_neighbourhood\_embedding} uses the neighbourhood
$(\phi^1_H)^{-1}(\Bo(C))$ and holds for every $n\ge2$.

\subsection*{What is verified}

The kernel checks the theorems of \texttt{Statement.lean}, \texttt{Computability.lean} and
\texttt{EpochComputability.lean} and all their dependencies in the development, the vendored library and Mathlib. The
trusted base is the kernel, the definitions in \texttt{Problem.lean} and the axioms \texttt{propext},
\texttt{Classical.choice} and \texttt{Quot.sound}; a command fails the build if the transitive axiom set of any of
these theorems contains anything else, in particular \texttt{sorryAx}.

Since Mathlib's real numbers use classical choice, axioms cannot distinguish a formula from an arbitrary existential
witness. Explicitness (level L1) and uniformity (level L3) are therefore checked by two metaprograms, which are not
themselves verified. \texttt{\#check\_explicit} traverses the non-proof dependency closure of a definition, through
project, vendored and Mathlib definitions, and accepts case distinctions on propositions, least-element searches over
integers, rationals and finite data with a non-classical decision procedure, uses of choice inside an explicit allow-list of Mathlib definitions
(\texttt{SymplecticPacking/Explicit/Whitelist.lean}), and definite descriptions after unique existence over types without a discrete factor. The $26$
definite descriptions reached from \texttt{General.packingHamiltonian} (inverses, roots and critical points, linear
solves, matrix square roots, solutions of ordinary differential equations, extrema and measures) are listed in Appendix~\Snum{sec:xpl}. The command is applied to \texttt{General.packingHamiltonian}, the Epoch Hamiltonians, the support
radius, the case embeddings and the algorithm terms. \texttt{\#check\_constructive} is stricter: the computational
closure may contain no classical constant at all, including classical decidability. It is applied to the two oracle
machines, so each is one fixed program in its oracle streams and cannot branch on an undecidable property of the
reals they name. \texttt{test/ExplicitTest.lean} contains terms both commands must reject, among them a search for the
least index whose enumerated Hamiltonian satisfies the specification, the same search wrapped as a unique pair in
$\N\times\R$, and a machine branching on $x<\frac13$ for a real $x$.

\emph{Independence.} The file \texttt{verification/Audit/Independence.lean} runs the command
\texttt{\#assert\_construction\_independence}, which traverses types, definitions and proof bodies of the packing, stability,
flow, corollary and algorithm theorems and fails if they reach the vendored packing predicate or existence theorems, the
holomorphic-line and two-ball obstruction proofs, or the curve-analysis modules of the vendored library. The converse
in Corollary~\ref{thm:main-sy} is checked separately: it may use the classical obstructions, but no sufficiency
theorem. A positive control detects its dependence on \texttt{adapted\_affine\_line}, and a regression test detects a
two-ball obstruction hidden inside a wrapper. The vendored declarations used are listed by
\texttt{verification/Audit/VendoredUsage.lean} (Appendix~\Snum{sec:xpl-vendored}): explicit charts, forms, primitives, Moser
fields and model coefficients.

Further statements carry source tags naming related Lean declarations; \texttt{verification/PAPER-INDEX.md} indexes
them, and the build checks that every tagged name exists and depends only on the three axioms.
The visible label \emph{full statement verified} denotes the full mathematical statement, up to the presentation
differences described above. The label \emph{needed consequence verified} denotes only the version used by the
construction, not every assertion or the entire domain of the printed statement. Appendix~\Snum{sec:xpl-interface}
records these restrictions, including return identities along smooth unit-vector families, transfer properties along
the constructed embeddings, and the coarser Hamiltonian support bound. An unqualified source tag is a reference,
without a full-coverage claim.
The scope labels are editorial correspondences for readers to review; the kernel checks the Lean statements.
The main correctness theorems concern the defined Hamiltonian and depend on every step of the construction.

\subsection*{Reproduction}

With Elan, Node.js and Tectonic installed, \texttt{lake exe cache get} followed by \texttt{bash scripts/verify.sh}
builds or incrementally checks the development with all axiom, explicitness, uniformity and independence gates, runs the
regression tests and the paper-tag check, and compiles both papers; every stage is fatal. The development is released
under Apache-2.0; the vendored library keeps its licence and notices.
The repository includes a short reviewer guide, \texttt{REVIEW.md}, and an optional project cache.
After building, \texttt{bash scripts/check\_kernel.sh} replays the complete proof closure of
\texttt{verification/Kernel.lean} into a fresh environment using the pinned toolchain's
\texttt{leanchecker}; this checks stored proof terms with the same Lean kernel.
A source rebuild from an empty project build directory checks elaboration without the project cache.

\endgroup
